\section{Bounded Booleanization, search, and finite quantification}
\label{sec:booleanization}

The single most reusable compilation step is to replace a bounded formula by a primitive-recursive Boolean whose value is proved equivalent to that formula.  This simultaneously solves three problems: bounded quantifiers disappear from induction invariants, finite case distinctions become intuitionistically legitimate, and bounded witnesses can be selected by least search.

\subsection{Boolean conventions}

Represent false by \(0\) and true by \(1=S(0)\).  Define primitive-recursive functions
\[
 \mathsf{not}(b),\quad \mathsf{and}(b,c),\quad \mathsf{or}(b,c),
 \quad \mathsf{isZero}(n),\quad \mathsf{eq}(m,n),\quad \mathsf{lt}(m,n),
\]
with outputs in \(\{0,1\}\).  The exact choice of definitions is unimportant provided their equations and range properties are proved in iPRA.  For example, equality can be computed by simultaneous bounded subtraction or by recursion on the two arguments; order can be computed from truncated subtraction.

For a Boolean term \(b\), write \([b]\) for the object formula \(b=1\).  Each Boolean connective has an object-iPRA reflection law, such as
\begin{align*}
 [\mathsf{not}(b)] &\leftrightarrow \neg[b],\\
 [\mathsf{and}(b,c)] &\leftrightarrow [b]\wedge[c],\\
 [\mathsf{or}(b,c)] &\leftrightarrow [b]\vee[c].
\end{align*}
The proof is a finite case distinction on the Boolean range of the arguments.  The range statement itself is quantifier-free and follows by primitive-recursive induction on the defining computation.

\subsection{Bounded formulas}

A bounded formula is generated from atomic equalities and inequalities by \(\wedge,\vee,\to,\neg\), and bounded quantifiers
\[
 \forall i<t\ A(i),\qquad \exists i<t\ A(i),
\]
where all bounds and terms are primitive recursive.  The formula may have arbitrary free numerical parameters, but the syntactic nesting depth is fixed.

\begin{theorem}[Structural Booleanization]
\label{thm:booleanization}
For every fixed bounded formula \(A(\vec{x})\), one can construct a primitive-recursive Boolean term \(\chi_A(\vec{x})\) and an object-iPRA derivation
\[
  \vdash_{\PRAAx} [\chi_A(\vec{x})]\leftrightarrow A(\vec{x}).
\]
Moreover, iPRA proves \(\chi_A(\vec{x})\le 1\).
\end{theorem}

\begin{proof}
Proceed by structural induction on \(A\).

For an equality \(s=t\), take \(\chi_A=\mathsf{eq}(s,t)\).  The defining correctness proof for \(\mathsf{eq}\) gives the equivalence.  Other decidable atomic relations are treated similarly.

For conjunction, disjunction, and negation, apply the corresponding Boolean operation and compose the induction hypotheses with its reflection law.  For implication, set
\[
  \chi_{A\to B}=\mathsf{or}(\mathsf{not}(\chi_A),\chi_B).
\]
Because the induction hypothesis supplies decidability of \(A\), intuitionistic logic proves
\(
eg A\vee B\leftrightarrow(A\to B)\).  Explicitly, the forward direction is immediate.  For the reverse direction, case-split on \(A\vee\neg A\); if \(A\), use the implication to obtain \(B\), and if \(\neg A\), take the left disjunct.  The case split is constructive because it is obtained from the Boolean value \(\chi_A\), not from unrestricted excluded middle.

For bounded universal quantification, define a fold
\begin{align*}
 \mathsf{All}_A(0,\vec{x}) &=1,\\
 \mathsf{All}_A(n+1,\vec{x})
  &=\mathsf{and}(\mathsf{All}_A(n,\vec{x}),\chi_A(n,\vec{x})).
\end{align*}
Use quantifier-free induction on \(n\) to prove
\[
 [\mathsf{All}_A(n,\vec{x})]
 \leftrightarrow \forall i<n\ A(i,\vec{x}).
\]
The right side in this displayed metatheoretic statement need not be the induction matrix.  In the actual object proof, the invariant is a quantifier-free equality relating the fold accumulator to the recursively computed Boolean table, or equivalently a Booleanized bounded universal.  The induction step uses only the recursion equation and the already proved reflection law for \(\chi_A\).  Specialize at the bound term \(t(\vec{x})\).

For bounded existential quantification, define
\begin{align*}
 \mathsf{Any}_A(0,\vec{x}) &=0,\\
 \mathsf{Any}_A(n+1,\vec{x})
  &=\mathsf{or}(\mathsf{Any}_A(n,\vec{x}),\chi_A(n,\vec{x})),
\end{align*}
and prove the analogous reflection theorem.  The output-range property is preserved by every Boolean constructor and fold.
\end{proof}

\subsection{A precise quantifier-free induction pattern}

It is worth spelling out how the universal-fold proof avoids placing a quantifier in the induction matrix.  Define an encoded prefix table
\[
 T_A(n,\vec{x})=\langle \chi_A(0,\vec{x}),\ldots,
                    \chi_A(n-1,\vec{x})\rangle
\]
by primitive recursion.  Define \(\mathsf{allTable}(T,n)\) by a bounded fold.  The induction invariant is the literal equation
\[
  a_n=\mathsf{allTable}(T_A(n,\vec{x}),n),
\]
where \(a_n\) is the accumulator generated by the recursion.  It is quantifier-free.  After the recursion theorem has established this equality, a separate finite-fold reflection lemma proves that \(a_n=1\) is equivalent to the bounded universal formula.  The proof therefore never uses induction on a quantified assertion.

The same pattern applies to finite sums, products, maxima, list membership, matrix identities over bounded indices, path checking, and coverage of a finite net.  One should regard a quantified mathematical statement as the \emph{specification} of a Boolean fold, not as the invariant itself.

\subsection{Least bounded search}

Let \(A(i,\vec{x})\) be bounded-decidable with Boolean \(\chi_A\).  Define a search state \((f,w)\), where \(f\in\{0,1\}\) records whether a witness has been found and \(w\) stores the first witness.  Initialize \((0,0)\).  At stage \(i\):
\[
 (f',w')=
 \begin{cases}
 (1,i),& f=0\text{ and }\chi_A(i,\vec{x})=1,\\
 (f,w),&\text{otherwise.}
 \end{cases}
\]
This is primitive recursive.

\begin{theorem}[Least bounded search]
\label{thm:least-search}
There is a primitive-recursive function \(\mu_{<B}A(\vec{x})\) such that iPRA proves:

\begin{enumerate}
\item if \(\exists i<B\ A(i,\vec{x})\), then
\(\mu_{<B}A<B\) and \(A(\mu_{<B}A,\vec{x})\);
\item if \(j<\mu_{<B}A\), then \(\neg A(j,\vec{x})\);
\item if no witness exists, the function returns a fixed default, say \(0\).
\end{enumerate}
\end{theorem}

\begin{proof}
Run the search state for \(B\) steps.  Use the quantifier-free invariant consisting of the following Boolean equalities:

\begin{itemize}
\item \(f\le 1\);
\item if \(f=1\), then \(w<i\) and \(\chi_A(w)=1\);
\item the Boolean fold \(\mathsf{Any}_A(i)\) equals \(f\);
\item a bounded Boolean fold confirms that every \(j<w\) has \(\chi_A(j)=0\).
\end{itemize}

The base case follows from the empty-fold equations.  The step is a finite Boolean calculation.  At time \(B\), nonemptiness implies \(\mathsf{Any}_A(B)=1\), hence \(f=1\), and the remaining invariant clauses yield validity and minimality.  Define \(\mu_{<B}A\) to be \(w\) when \(f=1\) and \(0\) otherwise.
\end{proof}

\subsection{Finite choice and its exact scope}

Bounded search yields a form of finite choice strong enough for most PFIG nodes.

\begin{corollary}[Bounded finite choice]
Suppose iPRA proves
\[
  \forall x<C\ \exists y<B(x)\ R(x,y)
\]
where \(R\) is bounded-decidable.  Then a primitive-recursive function \(f(x)\) can be defined by least bounded search, and iPRA proves \(R(x,f(x))\) for every \(x<C\).
\end{corollary}

There is no conflict with constructive restrictions: the domain and every search range are explicitly finite and decidable.  What does \emph{not} follow is a coherent selector from an infinite sequence of unrelated finite nonempty sets.  Coherence across precision is a separate theorem requiring compatibility and shrinking hypotheses, treated in \cref{sec:shrinking}.

\subsection{Booleanized certificates}

A certificate checker is a primitive-recursive Boolean \(\mathsf{check}(c,\vec{x})\).  The reflection theorem has the form
\[
 [\mathsf{check}(c,\vec{x})]\to \mathsf{Specification}(c,\vec{x}).
\]
The specification may contain quantifiers; only the checker and loop invariant must be quantifier-free.  This separation is useful for algebraic identities, paths in finite graphs, polynomial factorizations, triangulations, covering lists, determinant expansions, and finite analytic estimates.

A proof-producing checker node therefore needs two object derivations:

\begin{enumerate}
\item the finite computation establishes \(\mathsf{check}=1\);
\item reflection turns \(\mathsf{check}=1\) into the mathematical postcondition.
\end{enumerate}

The second proof is often theorem-specific but structurally simple.  It unfolds the certificate format and uses arithmetic lemmas; it does not repeat the search or trace argument.

\subsection{Compilation of finite sums and products}

For later applications, let \(a(i,\vec{x})\) be a primitive-recursive integer or rational-code expression.  Define
\begin{align*}
 \mathsf{Sum}(0)&=0,&
 \mathsf{Sum}(n+1)&=\mathsf{Sum}(n)+a(n),\\
 \mathsf{Prod}(0)&=1,&
 \mathsf{Prod}(n+1)&=\mathsf{Prod}(n)\cdot a(n).
\end{align*}
Their algebraic laws are proved by quantifier-free induction on equality invariants.  A nested finite sum is just a nested primitive recursion.  Uniform dimension does not create higher-type difficulty: matrices are bounded tables, matrix multiplication is a bounded sum, and determinant certificates can be checked by a bounded permutation sum or a division-free recurrence.

\subsection{Bounded minimization and finite nets}

If \(v(i)\) is a rational code for \(i<N\), compare values by exact cross multiplication and run a least-argmin fold.  iPRA proves that the selected index is within range and no candidate has smaller value.  If a finite net approximates a compact presentation with modulus \(m(p)\), then approximate minimization follows by:

\begin{enumerate}
\item compute the finite \(2^{-m(p)}\)-net;
\item evaluate a rational approximation to the objective at every net point;
\item choose the least approximate minimizer;
\item use the supplied continuity modulus to transfer the finite comparison to the represented object.
\end{enumerate}

The compactness content is entirely in the finite-net producer and its coverage certificate.  The minimization compiler is generic.

\subsection{Interface theorem for bounded logic}

Combining the preceding results gives the reusable statement needed by the graph compiler.

\begin{theorem}[Bounded-logic interface]
\label{thm:bounded-interface}
For every fixed bounded formula family and every fixed finite collection of folds, searches, sums, products, maxima, or table checks built from it, there are fixed primitive-recursive terms and object-iPRA derivations of their reflection, range, and least-witness laws.  All loop invariants used to establish these laws are quantifier-free.
\end{theorem}

This theorem is schematic in the paper and should become a syntax-directed Lean compiler.  Its syntax is intentionally limited to bounded formulas; unbounded quantifiers remain in specifications and are never evaluated by the compiler.
